\documentclass[tikz,border=8pt]{standalone}
\usepackage{tikz}
\usepackage{amsmath}
\usepackage{amssymb}
\usepackage{pgfplots}
\pgfplotsset{compat=1.18}
\usepackage{ctex}
\usetikzlibrary{calc,positioning,arrows.meta,matrix}

% LC 300 最长上升子序列（LIS）
% 数组 [10,9,2,5,3,7,101,18]
% dp[i] = 以 nums[i] 结尾的 LIS 长度
% 最终 LIS=4

\begin{document}
\begin{tikzpicture}[
  every node/.style={font=\small},
  cell/.style={draw=gray!70, minimum width=1.0cm, minimum height=0.8cm,
               inner sep=0pt, thick, fill=white},
  dpcell/.style={draw=gray!70, minimum width=1.0cm, minimum height=0.8cm,
                 inner sep=0pt, thick, fill=white},
  hicell/.style={draw=blue!60, minimum width=1.0cm, minimum height=0.8cm,
                 inner sep=0pt, thick, fill=blue!8},
  rescell/.style={draw=orange!70, minimum width=1.0cm, minimum height=0.8cm,
                  inner sep=0pt, very thick, fill=orange!8},
  xarr/.style={->,>=Stealth,thick},
]

%% ── 标题 ──────────────────────────────────────────────────
\node[font=\large\bfseries] at (4.5, 0.9)
  {最长上升子序列 LIS（LC 300）};
\node[font=\small, gray] at (4.5, 0.3)
  {数组 $[10,\,9,\,2,\,5,\,3,\,7,\,101,\,18]$，dp[i] = 以 nums[i] 结尾的 LIS 长度};

%% ── nums 数组行 ──────────────────────────────────────────
% 各位置 x = i*1.2 + 0.6 (i=0..7)
\def\xs{0.6}   % x start

\node[font=\scriptsize, gray] at (-0.7, -0.7) {索引};
\node[font=\scriptsize, gray] at (-0.7, -1.5) {nums};
\node[font=\scriptsize, gray] at (-0.7, -2.3) {dp};

\foreach \i/\v/\d/\style in
  {0/10/1/hicell, 1/9/1/dpcell, 2/2/1/dpcell,
   3/5/2/dpcell, 4/3/2/dpcell,
   5/7/3/dpcell, 6/101/4/rescell, 7/18/4/rescell} {
  \node[font=\scriptsize, gray] at (\i*1.2+0.6, -0.7) {\i};
  \node[cell] at (\i*1.2+0.6, -1.5) {\v};
  \node[\style] at (\i*1.2+0.6, -2.3) {\d};
}

%% ── 转移箭头（在数组下方，弧形） ──────────────────────────
% 箭头从 dp[j] → dp[i]，当 nums[j] < nums[i]
% 用 bend 弧形；只标关键路径

% dp[3]=2 来自 dp[2] (nums[2]=2 < nums[3]=5)
\draw[xarr, blue!60]
  (2*1.2+0.6, -2.75) to[out=-70, in=-110] (3*1.2+0.6, -2.75);
\node[font=\tiny, blue!60, below] at (2.85, -3.35) {2$<$5};

% dp[4]=2 来自 dp[2] (nums[2]=2 < nums[4]=3)
\draw[xarr, blue!60]
  (2*1.2+0.6, -2.75) to[out=-65, in=-115] (4*1.2+0.6, -2.75);
\node[font=\tiny, gray, below] at (3.5, -3.7) {2$<$3};

% dp[5]=3 来自 dp[3] (nums[3]=5 < nums[5]=7)
\draw[xarr, blue!70, thick]
  (3*1.2+0.6, -2.75) to[out=-60, in=-120] (5*1.2+0.6, -2.75);
\node[font=\tiny, blue!70, below] at (5.1, -3.5) {5$<$7};

% dp[5]=3 也可来自 dp[4] (nums[4]=3 < 7)
\draw[xarr, gray!50]
  (4*1.2+0.6, -2.75) to[out=-60, in=-120] (5*1.2+0.6, -2.75);

% dp[6]=4 来自 dp[5] (nums[5]=7 < nums[6]=101)
\draw[xarr, orange!70, very thick]
  (5*1.2+0.6, -2.75) to[out=-55, in=-125] (6*1.2+0.6, -2.75);
\node[font=\tiny, orange!70, below] at (7.05, -3.55) {7$<$101};

% dp[7]=4 来自 dp[5] (nums[5]=7 < 18)
\draw[xarr, orange!70, very thick]
  (5*1.2+0.6, -2.75) to[out=-50, in=-130] (7*1.2+0.6, -2.75);
\node[font=\tiny, orange!70, below] at (8.0, -3.85) {7$<$18};

%% ── 逐步演化说明 ──────────────────────────────────────────
\node[font=\normalsize\bfseries] at (4.5, -5.0) {dp 逐步计算过程};

% 表头
\node[font=\scriptsize, gray, anchor=west] at (0.0, -5.5) {步骤};
\node[font=\scriptsize, gray, anchor=west] at (1.2, -5.5) {i};
\node[font=\scriptsize, gray, anchor=west] at (2.2, -5.5) {nums[i]};
\node[font=\scriptsize, gray, anchor=west] at (4.0, -5.5) {来自 j（满足 nums[j]$<$nums[i]）};
\node[font=\scriptsize, gray, anchor=west] at (8.5, -5.5) {dp[i]};

\draw[gray!40] (0.0, -5.7) -- (10.5, -5.7);

\foreach \step/\ii/\nv/\from/\dpv in {
  1/0/10/{—（初始化）}/1,
  2/1/9/{—（无满足的 j）}/1,
  3/2/2/{—（无满足的 j）}/1,
  4/3/5/{j=2: dp[2]=1, 2<5 → dp[3]=2}/2,
  5/4/3/{j=2: dp[2]=1, 2<3 → dp[4]=2}/2,
  6/5/7/{j=3: dp[3]=2, 5<7 → dp[5]=3}/3,
  7/6/101/{j=5: dp[5]=3, 7<101 → dp[6]=4}/4,
  8/7/18/{j=5: dp[5]=3, 7<18 → dp[7]=4}/4%
} {
  \pgfmathtruncatemacro{\yy}{-6 - (\step-1)*0.55}
  \node[font=\scriptsize, anchor=west] at (0.0, \yy) {\step};
  \node[font=\scriptsize, anchor=west] at (1.2, \yy) {\ii};
  \node[font=\scriptsize, anchor=west] at (2.2, \yy) {\nv};
  \node[font=\scriptsize, anchor=west] at (4.0, \yy) {\from};
  \node[font=\scriptsize\bfseries, anchor=west] at (8.5, \yy) {\dpv};
}

\draw[gray!40] (0.0, -10.55) -- (10.5, -10.55);

%% ── 最终 LIS 路径展示 ──────────────────────────────────────
\node[draw=gray!50, fill=white, rounded corners=3pt, font=\small,
      inner sep=6pt, align=left]
  at (4.5, -11.5)
  {\textbf{最终 LIS = 4}，一条最优路径：$2 \to 5 \to 7 \to 101$（或 $2 \to 3 \to 7 \to 18$）\\
   \textbf{算法复杂度：}$O(n^2)$ DP \quad 优化：$O(n \log n)$ patience sorting};

\end{tikzpicture}
\end{document}
